\section*{Bab \ref{ch:b1:matrices} --- Matriks}

\begin{solution}{exo:b1:matrices:1}
\[
AB = \begin{pmatrix} 2 & 1\\ 1 & 0\end{pmatrix},
\quad
BA = \begin{pmatrix} 0 & 1\\ 1 & 2\end{pmatrix},
\quad
A^2 - B^2 = \begin{pmatrix} 1 & 4\\ 0 & 1\end{pmatrix} - I
= \begin{pmatrix} 0 & 4\\ 0 & 0\end{pmatrix},
\]
\[
(A+B)(A-B) = A^2 - AB + BA - B^2
= \begin{pmatrix} 0 & 4\\ 0 & 0\end{pmatrix} +
\begin{pmatrix} -2 & 0\\ 0 & 2 \end{pmatrix}
= \begin{pmatrix} -2 & 4\\ 0 & 2\end{pmatrix}.
\]
Keduanya berbeda sebesar $BA - AB \neq 0$: karena kesamaan $(a+b)(a-b) = a^2 -
b^2$ menuntut kekomutatifan, yang gagal di sini.
\end{solution}

\begin{solution}{exo:b1:matrices:2}
$M$: $L_2 \leftarrow L_2 - 2L_1$ membunuh baris keduanya; $L_3
\leftarrow L_3 - L_1$ memberikan $(0, -1, -2)$. Dua poros:
$\operatorname{rk} M = 2$.

$N$: $L_3 \leftarrow L_3 - L_1$ memberikan $(0,1,1,1) = L_2$; lalu $L_3
\leftarrow L_3 - L_2 = 0$. Dua poros: $\operatorname{rk} N = 2$.
\end{solution}

\begin{solution}{exo:b1:matrices:3}
Mereduksi $(A \mid I_3)$: $L_2 \leftarrow L_2 - 2L_1$, $L_3
\leftarrow L_3 - L_1$:
\[
\begin{pmatrix}
1 & 0 & 1 & 1 & 0 & 0\\
0 & 1 & -1 & -2 & 1 & 0\\
0 & 1 & 0 & -1 & 0 & 1
\end{pmatrix}
\xrightarrow{L_3 \leftarrow L_3 - L_2}
\begin{pmatrix}
1 & 0 & 1 & 1 & 0 & 0\\
0 & 1 & -1 & -2 & 1 & 0\\
0 & 0 & 1 & 1 & -1 & 1
\end{pmatrix},
\]
lalu $L_1 \leftarrow L_1 - L_3$, $L_2 \leftarrow L_2 + L_3$:
\[
A^{-1} = \begin{pmatrix}
0 & 1 & -1\\
-1 & 0 & 1\\
1 & -1 & 1
\end{pmatrix}.
\]
\emph{Periksa:} baris pertama $A$ kali kolom pertama $A^{-1}$: $1
\cdot 0 + 0\cdot(-1) + 1\cdot 1 = 1$; kali kolom kedua: $1 - 0 -
1 = 0$; kali yang ketiga: $-1 + 0 + 1 = 0$.
\end{solution}

\begin{solution}{exo:b1:matrices:4}
$u(1) = 1$, $u(X) = X + 1$, $u(X^2) = X^2 + 2X + 1$: kolom
\omterm{prop:b1:vspaces:coordinates}{koordinat} pada $(1, X, X^2)$ memberikan
\[
M = \begin{pmatrix}
1 & 1 & 1\\
0 & 1 & 2\\
0 & 0 & 1
\end{pmatrix}.
\]
$u$ dapat dibalik karena ia mempunyai invers yang jelas $P \mapsto P(X
- 1)$ (yaitu komposisi penyulihan). Matriksnya diperoleh dengan
cara yang sama dari $u^{-1}(X^k) = (X-1)^k$:
\[
M^{-1} = \begin{pmatrix}
1 & -1 & 1\\
0 & 1 & -2\\
0 & 0 & 1
\end{pmatrix}.
\]
\end{solution}

\begin{solution}{exo:b1:matrices:5}
$N = \begin{pmatrix} 0 & 1\\ 0 & 0\end{pmatrix}$, $N^2 = 0$. Karena
$2I$ dan $N$ berkomutasi, penjabaran binomialnya terpenggal setelah dua
suku:
\[
A^k = (2I + N)^k = 2^k I + k\,2^{k-1} N
= \begin{pmatrix} 2^k & k\,2^{k-1}\\ 0 & 2^k\end{pmatrix}.
\]
(Periksa $k = 2$: $A^2 = \begin{pmatrix}4 & 4\\ 0 & 4\end{pmatrix}$,
yang benar lewat hasil kali langsung.)
\end{solution}

\begin{solution}{exo:b1:matrices:6}
\omterm{def:b1:matrices:transpose}{Trace}: $\operatorname{tr}(AB - BA) = \operatorname{tr}(AB) -
\operatorname{tr}(BA) = 0$ (\cref{def:b1:matrices:transpose}),
sedangkan $\operatorname{tr}(I_n) = n \neq 0$ di dalam $\R$ atau $\C$. Jadi tak ada
penyelesaian. (Pada ruang berdimensi takhingga kesamaannya \emph{memang}
terwujudkan --- penurunan dan perkalian dengan $x$ memenuhinya
--- justru karena tak ada \omterm{def:b1:matrices:transpose}{trace} di sana.)
\end{solution}

\begin{solution}{exo:b1:matrices:7}
Hasil kali teleskopik, dengan semua pangkat $A$ yang berkomutasi:
\[
(I - A)(I + A + \dots + A^{m-1}) = I - A^m = I ,
\]
dan \cref{prop:b1:matrices:ring} meningkatkan invers sepihaknya. Adapun untuk
penerapannya: matriks yang diberikan adalah $I + N$ dengan
\[
N = \begin{pmatrix} 0 & 2 & 3\\ 0 & 0 & 2\\ 0&0&0 \end{pmatrix},
\quad
N^2 = \begin{pmatrix} 0&0&4\\ 0&0&0\\ 0&0&0\end{pmatrix},
\quad N^3 = 0 ,
\]
sehingga, dengan mengganti $A$ dengan $-N$ pada rumusnya:
\[
(I + N)^{-1} = I - N + N^2 =
\begin{pmatrix}
1 & -2 & 1\\
0 & 1 & -2\\
0 & 0 & 1
\end{pmatrix}.
\]
\end{solution}

\begin{solution}{exo:b1:matrices:8}
$A^2 = A$: jadi endomorfisma $a$ merupakan \omterm{def:b1:linmaps:projection}{proyeksi}
(\cref{thm:b1:linmaps:projchar}), $E = \operatorname{im} a \oplus
\ker a$ dengan $\dim\operatorname{im} a = r = \operatorname{rk} A$. Pada
\omterm{def:b1:vspaces:free}{basis} yang disesuaikan dengan penguraian ini ($r$ vektor peta,
lalu \omterm{def:b1:vspaces:free}{basis} kernelnya), matriks $a$ adalah
$\begin{pmatrix} I_r & 0\\ 0 & 0\end{pmatrix}$, yang bertrace $r$. Adapun
\omterm{def:b1:matrices:transpose}{tracenya} invarian terhadap \omterm{def:b1:matrices:changeofbasis}{perubahan basis}:
$\operatorname{tr}(P^{-1}MP) = \operatorname{tr}(MPP^{-1}) =
\operatorname{tr} M$ menurut kesamaan berdaurnya. Maka
$\operatorname{tr} A = r = \operatorname{rk} A$.
\end{solution}

\begin{solution}{exo:b1:matrices:9}
$J^2 = nJ$ (karena setiap entri $J^2$ menjumlahkan $n$ satuan). Carilah $M^{-1} =
\alpha I + \beta J$:
\[
(aI + bJ)(\alpha I + \beta J)
= a\alpha\, I + (a\beta + b\alpha + nb\beta)\, J .
\]
Ini sama dengan $I$ jika dan hanya jika $a\alpha = 1$ dan $a\beta + b\alpha + nb\beta =
0$, yaitu $\alpha = \frac1a$ dan $\beta(a + nb) = -\frac ba$. Jika $a
\neq 0$ dan $a + nb \neq 0$:
\[
M^{-1} = \frac 1a I - \frac{b}{a(a + nb)}\, J .
\]
Sebaliknya, jika $a = 0$: maka $M = bJ$ ber-rank $\leq 1 < n$ (untuk $n \geq
2$): jadi tak dapat dibalik ($n = 1$ merupakan kasus skalarnya). Jika $a + nb = 0$:
maka vektor $v = (1, \dots, 1)^{\mathsf T}$ memenuhi $Mv = (a +
nb)v = 0$ dengan $v \neq 0$: jadi tak dapat dibalik. Sehingga $M \in GL_n \iff a
\neq 0$ dan $a + nb \neq 0$.
\end{solution}

\begin{solution}{exo:b1:matrices:10}
\emph{Jumlah:} $\operatorname{im}(A + B) \subseteq \operatorname{im} A
+ \operatorname{im} B$ (karena setiap $(A+B)x = Ax + Bx$), dan Grassmann
membatasi dimensi sebuah jumlah oleh jumlah dimensinya.

\emph{Sylvester:} misalkan $a$ \omterm{def:b1:logic:map}{pemetaan} $A$ yang dibatasi ke $V =
\operatorname{im} B$ (yang berdimensi $\operatorname{rk} B$). Petanya adalah
$\operatorname{im}(AB)$ (karena $a(Bx) = ABx$), dan rank--nulitas di dalam $V$:
\[
\operatorname{rk} B = \dim\ker(a_{|V}) + \operatorname{rk}(AB) .
\]
Kini $\ker(a_{|V}) \subseteq \ker A$, yang berdimensi $n -
\operatorname{rk} A$: sehingga
\[
\operatorname{rk}(AB) \geq \operatorname{rk} B - (n -
\operatorname{rk} A) = \operatorname{rk} A + \operatorname{rk} B -
n . \qedhere
\]
\end{solution}

\begin{solution}{exo:b1:matrices:11}
\begin{enumerate}
  \item Entri demi entri, $(AD)_{ij} = a_{ij}\,d_j$ dan $(DA)_{ij} =
        d_i\,a_{ij}$. Jadi $AD = DA$ jika dan hanya jika $a_{ij}(d_j - d_i) = 0$
        untuk semua $i, j$; dan ketika $i \neq j$ faktor $d_j - d_i$
        taknol, yang memaksa $a_{ij} = 0$: sehingga $A$ diagonal.
        Sebaliknya matriks diagonal saling berkomutasi.
  \item Jika $A$ berkomutasi dengan setiap matriks, maka ia berkomutasi dengan
        $\operatorname{diag}(1, 2, \dots, n)$, sehingga $A =
        \operatorname{diag}(\lambda_1, \dots, \lambda_n)$ menurut (1).
        Lalu $A E_{ij} = \lambda_i E_{ij}$ (karena hanya baris $i$ pada
        $E_{ij}$ yang bertahan) sedangkan $E_{ij} A = \lambda_j E_{ij}$:
        jadi berkomutasi dengan $E_{ij}$ memaksa $\lambda_i = \lambda_j$.
        Maka $A = \lambda I_n$; dan matriks skalar memang berkomutasi
        dengan segalanya. Jadi pusat $\mathcal{M}_n(K)$ adalah
        $K\,I_n$.
\end{enumerate}
\end{solution}

\begin{solution}{exo:b1:matrices:12}
\begin{enumerate}
  \item Jika $\operatorname{rk} A = 1$: maka peta $A$ berupa garis
        $\operatorname{Vect}(C)$, $C \neq 0$, sehingga kolom ke-$j$
        pada $A$ adalah $\ell_j\,C$ bagi skalar $\ell_j$ (yang tak
        semuanya nol), yaitu $A = C L$ dengan $L = (\ell_1, \dots,
        \ell_n) \neq 0$. Sebaliknya jika $A = CL \neq 0$, maka semua
        kolomnya kelipatan $C$: jadi ber-rank $1$.
  \item $A^2 = C\,(L C)\,L$, dan $LC$ merupakan skalar $\sum_i
        \ell_i c_i = \operatorname{tr}(CL) = \operatorname{tr}
        A$. Jadi $A^2 = (\operatorname{tr} A)\,A$, sehingga lewat
        induksi $A^m = (\operatorname{tr} A)^{m-1} A$. Jika
        $\operatorname{tr} A \neq 0$, maka tak ada pangkatnya yang lenyap; jika
        $\operatorname{tr} A = 0$, maka $A^2 = 0$: sehingga matriks
        ber-rank satu bersifat nilpoten jika dan hanya jika \omterm{def:b1:matrices:transpose}{tracenya} nol.
  \item Dengan $t = \operatorname{tr} A \neq -1$:
        \[
        (I_n + A)\Bigl(I_n - \frac{A}{1 + t}\Bigr)
        = I_n + A - \frac{A + A^2}{1 + t}
        = I_n + A - \frac{(1 + t)A}{1 + t} = I_n ,
        \]
        dengan memakai $A^2 = tA$. Jika $t = -1$: $(I_n + A)A = A + A^2 = A
        - A = 0$ dengan $A \neq 0$, sehingga $I_n + A$ membunuh setiap
        kolom (taknol) $A$: jadi tak \omterm{def:b1:logic:inj}{injektif}, tak dapat dibalik.
\end{enumerate}
\end{solution}

\begin{solution}{pb:b1:matrices:1}
\textbf{1.} Pembagian Euclid atas $X^n$ oleh yang \omterm{def:b1:poly:def}{monik} berderajat $2$
yaitu $D$ (\cref{thm:b1:poly:division}): hasil bagi dan sisanya ada
dan tunggal, dan sisanya berderajat $\leq 1$: $X^n =
Q_n D + a_n X + b_n$. Untuk $n = 0$: $Q_0 = 0$, $(a_0, b_0) = (0,
1)$; untuk $n = 1$: $(a_1, b_1) = (1, 0)$.

\textbf{2.} Kalikanlah dengan $X$ lalu reduksikanlah $X^2 = D + sX - p$:
\[
X^{n+1} = X Q_n D + a_n X^2 + b_n X
= (X Q_n + a_n)\,D + (s\,a_n + b_n)\,X - p\,a_n .
\]
Ungkapan terakhirnya berbentuk sisa (yaitu berderajat $\leq 1$), sehingga menurut
ketunggalannya $a_{n+1} = s a_n + b_n$ dan $b_{n+1} = -p a_n$.
Menyulihkan $b_{n+1} = -pa_n$ ke dalam $a_{n+2} = s a_{n+1} +
b_{n+1}$ memberikan $a_{n+2} = s\,a_{n+1} - p\,a_n$.

\textbf{3.} Nilailah $X^n = Q_n D + a_n X + b_n$ pada akarnya:
$\lambda^n = a_n\lambda + b_n$ dan $\mu^n = a_n\mu + b_n$.
Dengan mengurangkan lalu \omterm{def:b1:arith:divides}{membaginya} dengan $\lambda - \mu \neq 0$:
\[
a_n = \frac{\lambda^n - \mu^n}{\lambda - \mu},
\qquad
b_n = \lambda^n - a_n\lambda
= \frac{\lambda\mu^n - \mu\lambda^n}{\lambda - \mu} .
\]

\textbf{4.} Pada akar rangkapnya: $\lambda^n = a_n\lambda + b_n$.
Menurunkan kesamaannya, $nX^{n-1} = Q_n'\,(X - \lambda)^2
+ 2Q_n\,(X - \lambda) + a_n$, lalu menilainya di $\lambda$: $a_n
= n\lambda^{n-1}$; lalu $b_n = \lambda^n - n\lambda^{n} = (1 -
n)\lambda^{n}$.

\textbf{5.} Untuk $P = \sum_i p_i X^i$ dan $Q = \sum_j q_j X^j$,
\[
P(M)\,Q(M) = \sum_{i,j} p_i q_j M^{i+j} = (PQ)(M),
\]
karena pangkat satu matriks tunggal $M$ saling berkomutasi
(sedangkan jumlahnya jelas lewat kelinearan). Jika $D(M) = 0$, maka menyulihkan $M$
ke dalam $X^n = Q_n D + a_n X + b_n$ memberikan $M^n = Q_n(M)\,D(M) + a_n
M + b_n I = a_n M + b_n I$.

\textbf{6.} Hasil kali langsungnya:
\[
A^2 = \begin{pmatrix}
a^2 + bc & b(a + d)\\
c(a + d) & d^2 + bc
\end{pmatrix},
\qquad
s A = \begin{pmatrix}
a(a+d) & b(a+d)\\
c(a+d) & d(a+d)
\end{pmatrix},
\]
sehingga $A^2 - sA$ berentri luar diagonal nol dan berentri diagonal
$a^2 + bc - a^2 - ad = bc - ad = -p$: jadi $A^2 - sA + pI_2 = 0$.

\textbf{7.} Dengan $A' = \begin{pmatrix} a' & b'\\ c' &
d'\end{pmatrix}$, menjabarkan $p(AA') = (aa' + bc')(cb' + dd') -
(ab' + bd')(ca' + dc')$: suku $aa'cb'$ dan $ab'ca'$ saling menghapus,
suku $bc'dd'$ dan $bd'dc'$ saling menghapus, dan yang tersisa adalah
\[
aa'dd' - bca'd' + bcb'c' - adb'c'
= (ad - bc)(a'd' - b'c') = p(A)\,p(A').
\]
Jika $p \neq 0$, maka Cayley--Hamilton memberikan $A\,\bigl(\tfrac1p(sI_2 -
A)\bigr) = \tfrac1p(sA - A^2) = I_2$, dan dari situlah inversnya (dan
\cref{prop:b1:matrices:ring} menjadikannya dua sisi). Jika $p = 0$ dan
$A$ dapat dibalik, maka kesifatan perkaliannya memberikan $1 = p(I_2) =
p(A)\,p(A^{-1}) = 0$: mustahil. Jadi $A \in GL_2 \iff p \neq 0$.

\textbf{8.} $s = 3$, $p = 2$, $D = X^2 - 3X + 2 = (X - 1)(X -
2)$: $\lambda = 2$, $\mu = 1$, sehingga $a_n = 2^n - 1$ dan $b_n = 2 -
2^n$ (pertanyaan 3). Maka
\[
A^n = (2^n - 1)A + (2 - 2^n)I
= \begin{pmatrix} 1 & 2^n - 1\\ 0 & 2^n \end{pmatrix}.
\]
Periksa: $A^2 = \begin{pmatrix} 1 & 3\\ 0 & 4\end{pmatrix}$ baik
lewat rumusnya maupun lewat pengkuadratan langsung.

\textbf{9.} $s = 4$, $p = 3\cdot1 - 1\cdot(-1) = 4$: $D = X^2 -
4X + 4 = (X - 2)^2$, dengan akar rangkap $\lambda = 2$. Pertanyaan 4: $a_n
= n\,2^{n-1}$, $b_n = (1 - n)2^n$, sehingga
\[
A^n = n\,2^{n-1}A + (1 - n)2^n I
= 2^{n-1}\begin{pmatrix} n + 2 & n\\ -n & 2 - n
\end{pmatrix}.
\]
Di $n = 2$: $2\begin{pmatrix} 4 & 2\\ -2 & 0\end{pmatrix} =
\begin{pmatrix} 8 & 4\\ -4 & 0 \end{pmatrix}$, yaitu $A^2$
yang dihitung langsung.

\textbf{10.} Induksi: $F^1 = \begin{pmatrix} F_2 & F_1\\ F_1 &
F_0\end{pmatrix}$, dan
\[
F^{n+1} = F^n F =
\begin{pmatrix} F_{n+1} + F_n & F_{n+1}\\
F_n + F_{n-1} & F_n \end{pmatrix}
= \begin{pmatrix} F_{n+2} & F_{n+1}\\ F_{n+1} & F_n
\end{pmatrix}.
\]
Di sini $s = 1$, $p = -1$, $D = X^2 - X - 1$ berakar $\varphi,
\psi$ ($\varphi - \psi = \sqrt5$, $\varphi\psi = -1$). Adapun
barisan $(F_n)$ mempunyai $F_0 = 0 = a_0$, $F_1 = 1 = a_1$ dan menuruti
rekurensi yang sama dengan $(a_n)$: sehingga $F_n = a_n = (\varphi^n -
\psi^n)/\sqrt5$, yaitu rumus Binet. Cassini: dengan menerapkan kesifatan
perkalian pertanyaan 7 pada $F^n$,
\[
F_{n+1}F_{n-1} - F_n^2 = p(F^n) = p(F)^n = (-1)^n .
\]

\textbf{11.} Syaratnya \omterm{def:b1:linmaps:def}{linear} dan memuat barisan
nol: jadi sebuah \omterm{def:b1:vspaces:subspace}{subruang}. Lewat induksi $u_0, u_1$ menentukan $u$
secara \omterm{def:b1:linmaps:def}{linear}, dan setiap pasangan nilai awalnya terwujud oleh
tepat satu penyelesaian: jadi seperti pada \cref{exo:b1:findim:10}, $E_D$
terparameterkan secara \omterm{def:b1:logic:inj}{bijektif} dan \omterm{def:b1:linmaps:def}{linear} oleh $(u_0, u_1) \in K^2$:
sehingga $\dim E_D = 2$.

\textbf{12.} $(a_n)$ menuruti rekurensinya (pertanyaan 2) dengan $a_0
= 0$, $a_1 = 1$. Demikian pula $(b_n)$: $b_{n+2} = -p\,a_{n+1} =
-p(s a_n + b_n) = s\,b_{n+1} - p\,b_n$ (dengan memakai $b_{n+1} = -pa_n$
dua kali), dengan $b_0 = 1$, $b_1 = 0$. Maka kombinasi $v_n =
u_1 a_n + u_0 b_n$ merupakan penyelesaian dengan $v_0 = u_0$, $v_1 =
u_1$; sedangkan dua penyelesaian bernilai awal sama berimpit
(lewat induksi), sehingga $u_n = u_1 a_n + u_0 b_n$ untuk semua $n$.

\textbf{13.} $(\lambda^n)$ merupakan penyelesaian jika dan hanya jika $\lambda^{n+2} =
s\lambda^{n+1} - p\lambda^n$ untuk semua $n$, yaitu $D(\lambda) = 0$
(setelah dibagi $\lambda^n \neq 0$; perhatikanlah $\lambda, \mu \neq
0$ karena $p = \lambda\mu \neq 0$). Adapun kebebasan
$\bigl((\lambda^n), (\mu^n)\bigr)$: sebuah hubungan di $n = 0, 1$
memberikan $c + c' = 0$, $c\lambda + c'\mu = 0$, sehingga $c(\lambda - \mu)
= 0$: jadi $c = c' = 0$. Dua vektor \omterm{def:b1:vspaces:free}{bebas} dalam dimensi $2$: yaitu \omterm{def:b1:vspaces:free}{basis}.
Akar rangkapnya: $\bigl((n\lambda^n)\bigr)$ merupakan penyelesaian karena,
dengan $s = 2\lambda$, $p = \lambda^2$:
\[
s(n+1)\lambda^{n+1} - p\,n\lambda^n
= \lambda^{n+2}\bigl(2(n+1) - n\bigr) = (n+2)\lambda^{n+2} ;
\]
kebebasannya di $n = 0, 1$: $c = 0$, lalu $c'\lambda = 0$ dengan
$\lambda \neq 0$.

\textbf{14.} $D = X^2 - X - 6 = (X - 3)(X + 2)$. Penyelesaian
umumnya $u_n = A\,3^n + B(-2)^n$; syarat awalnya memberikan
$A + B = 1$ dan $3A - 2B = 8$, sehingga $A = 2$, $B = -1$:
\[
u_n = 2\cdot 3^n - (-2)^n .
\]
Periksa: $u_2 = 18 - 4 = 14 = u_1 + 6u_0$; $u_3 = 54 + 8 = 62 =
u_2 + 6u_1 = 14 + 48$.

\textbf{15.} $C\begin{pmatrix} u_n\\ u_{n+1}\end{pmatrix} =
\begin{pmatrix} u_{n+1}\\ -p\,u_n + s\,u_{n+1}\end{pmatrix} =
\begin{pmatrix} u_{n+1}\\ u_{n+2}\end{pmatrix}$, dan induksi
memberikan rumusnya dengan $C^n$. Lebih-lebih $\operatorname{tr} C = 0
+ s = s$ dan $p(C) = 0\cdot s - 1\cdot(-p) = p$: sehingga matriks
pendampingnya mempunyai tepat $D$ sebagai \omterm{def:b1:poly:def}{polinomial} Cayley--Hamiltonnya.

\textbf{16.} Bagi sebarang penyelesaian $u$ atas $u_{n+3} = \alpha u_{n+2} +
\beta u_{n+1} + \gamma u_n$, vektor keadaannya $v_n = (u_n,
u_{n+1}, u_{n+2})^{\mathsf T}$ memenuhi $C_3 v_n = v_{n+1}$
(karena kedua baris pertamanya menggeser, sedangkan baris terakhirnya menerapkan
rekurensinya). Maka
\[
D_3(C_3)\,v_0 = v_3 - \alpha v_2 - \beta v_1 - \gamma v_0 ,
\]
yang ketiga komponennya adalah $u_{k+3} - \alpha u_{k+2} - \beta
u_{k+1} - \gamma u_k = 0$ ($k = 0, 1, 2$). Dan ketika keadaan awalnya
$v_0 = (u_0, u_1, u_2)^{\mathsf T}$ menjelajahi \emph{seluruh}
$K^3$ (karena nilai awalnya \omterm{def:b1:vspaces:free}{bebas}), matriks $D_3(C_3)$ membunuh
setiap vektor: jadi $D_3(C_3) = 0$.

\textbf{17.} Tulislah $X^n = Q\,D_3 + R_n$ dengan $\deg R_n \leq 2$
lalu nilailah pada setiap akarnya: $\lambda_i^n = R_n(\lambda_i)$. Jadi
$R_n$ merupakan \omterm{def:b1:poly:def}{polinomial} berderajat $\leq 2$ yang menginterpolasi ketiga
nilai $\lambda_i^n$ pada ketiga simpul berbeda $\lambda_i$: sehingga menurut
ketunggalan pada \cref{thm:b1:poly:lagrange}, $R_n = \sum_i
\lambda_i^n L_i$ dengan $(L_i)$ sebagai \omterm{def:b1:vspaces:free}{basis} Lagrange simpulnya.
Menyulihkan $C_3$ (pertanyaan 5 dan 16):
\[
C_3^{\,n} = R_n(C_3) = \sum_{i=1}^{3} \lambda_i^n\,L_i(C_3),
\]
dengan ketiga matriks $L_i(C_3)$ yang tak bergantung pada $n$: sehingga setiap
entri $C_3^{\,n}$ merupakan kombinasi tetap atas $\lambda_1^n,
\lambda_2^n, \lambda_3^n$.

\textbf{18.} $D_3 = X^3 - 2X^2 - X + 2 = (X-1)(X+1)(X-2)$.
Penyelesaian umumnya $u_n = A + B(-1)^n + C\,2^n$. Syarat
awalnya: $A + B + C = 0$, $A - B + 2C = 1$, $A + B + 4C = 1$.
Mengurangkan yang pertama dari yang ketiga: $3C = 1$, $C = \frac13$;
lalu $A + B = -\frac13$ dan $A - B = \frac13$: $A = 0$, $B =
-\frac13$. Maka
\[
u_n = \frac{2^n - (-1)^n}{3}
\]
(yaitu bilangan Jacobsthal). Periksa: $u_3 = \frac{8 + 1}{3} = 3 =
2u_2 + u_1 - 2u_0 = 2 + 1 - 0$.

\textbf{19.} Ekspansi Taylor atas \omterm{def:b1:poly:def}{polinomial} $X^n$ di
$\lambda$:
\[
X^n = \sum_{k=0}^{n} \binom nk \lambda^{n-k}(X - \lambda)^k ,
\]
dan semua sukunya dengan $k \geq 3$ terbagi oleh $(X -
\lambda)^3$: sehingga sisanya adalah
\[
R_n = \lambda^n + n\lambda^{n-1}(X - \lambda) + \binom
n2\lambda^{n-2}(X - \lambda)^2 .
\]
Untuk $M = \lambda I + N$ dengan $N^3 = 0$: $(M - \lambda I)^3 = N^3
= 0$, sehingga pertanyaan 5 memberikan
\[
M^n = \lambda^n I + n\lambda^{n-1} N + \binom n2
\lambda^{n-2} N^2 ,
\]
yaitu tepat penjabaran binomial atas $(\lambda I + N)^n$
yang terpenggal di $N^2$ --- jadi kedua metodenya sepakat.

\textbf{20.} Ruang penyelesaiannya berdimensi $3$ (dengan
parameterisasi $(u_0, u_1, u_2)$ yang sama seperti pada pertanyaan 11), dan
setiap $(\lambda_i^n)$ merupakan penyelesaian. Kebebasannya: andaikan
$c_1\lambda_1^n + c_2\lambda_2^n + c_3\lambda_3^n = 0$ untuk $n =
0, 1, 2$. Tetapkanlah $i$ lalu misalkan $L_i = \sum_{k \leq 2} p_k X^k$ sebagai
\omterm{def:b1:poly:def}{polinomial} Lagrange simpulnya dengan $L_i(\lambda_j) =
\delta_{ij}$. Maka
\[
0 = \sum_{k=0}^{2} p_k\Bigl(\sum_j c_j\lambda_j^k\Bigr)
= \sum_j c_j\,L_i(\lambda_j) = c_i .
\]
Jadi semua $c_i = 0$: yaitu tiga penyelesaian \omterm{def:b1:vspaces:free}{bebas} dalam dimensi $3$, sehingga sebuah
\omterm{def:b1:vspaces:free}{basis}; dan penyelesaian umumnya adalah $c_1\lambda_1^n +
c_2\lambda_2^n + c_3\lambda_3^n$.

\textbf{21.} Dari $F_k = F_{k+2} - F_{k+1}$, jumlahnya berteleskop:
\[
\sum_{k=1}^{n} F_k = \sum_{k=1}^{n}\bigl(F_{k+2} - F_{k+1}\bigr)
= F_{n+2} - F_2 = F_{n+2} - 1 .
\]

\textbf{22.} Ambillah entri $(1,2)$ pada $F^{m+n} = F^m F^n$: sisi
kirinya adalah $F_{m+n}$; sedangkan sisi kanannya (baris $1$ pada $F^m$)
kali (kolom $2$ pada $F^n$), yaitu $F_{m+1}F_n + F_m F_{n-1}$.
Dengan $m = n$:
\[
F_{2n} = F_{n+1}F_n + F_nF_{n-1} = F_n\,(F_{n+1} + F_{n-1}).
\]

\textbf{23.} Menurut Binet, $F_n - \dfrac{\varphi^n}{\sqrt5} =
-\dfrac{\psi^n}{\sqrt5}$, dan $\abs\psi = \frac{\sqrt5 - 1}2 <
1$, sehingga
\[
\Bigl|F_n - \frac{\varphi^n}{\sqrt5}\Bigr|
\leq \frac{1}{\sqrt5} < \frac12
\qquad (n \geq 0):
\]
$F_n$ merupakan bilangan bulat terdekat dengan $\varphi^n/\sqrt5$.

\textbf{24.} $t_n = F_{n+1} + F_{n-1}$ merupakan kombinasi atas
barisan Fibonacci yang tergeser, sehingga memenuhi
rekurensi yang sama: $t_{n+2} = t_{n+1} + t_n$; dan $t_1 = F_2 + F_0 =
1$, $t_2 = F_3 + F_1 = 3$: yaitu bilangan Lucas $L_n$.
Adapun barisan $\varphi^n + \psi^n$ merupakan penyelesaian bernilai
dua yang pertama sama ($\varphi + \psi = 1$, $\varphi^2 + \psi^2 = (
\varphi + \psi)^2 - 2\varphi\psi = 3$), sehingga $L_n = \varphi^n +
\psi^n$. Akhirnya
\[
F_n L_n = \frac{(\varphi^n - \psi^n)(\varphi^n +
\psi^n)}{\sqrt5} = \frac{\varphi^{2n} - \psi^{2n}}{\sqrt5} =
F_{2n},
\]
yang memulihkan pertanyaan 22.

\textbf{25.} (i) Kelima matriks $I, A, A^2, A^3, A^4$ hidup di
dalam ruang berdimensi $4$ yaitu $\mathcal{M}_2(K)$, sehingga \emph{ada}
\omterm{def:b1:poly:def}{polinomial} taknol berderajat $\leq 4$ yang membunuh $A$; lalu Bagian II mempertajamnya
menjadi kuadratik yang eksplisit $A^2 = sA - pI$, yang mengunci semua
pangkatnya ke dalam bidang $\operatorname{Vect}(I, A)$. (ii)
Pembagian Euclid mereduksi $X^n$ modulo kuadratik itu, dan
kedua koefisien sisanya menuruti rekurensi bersuku dua
$a_{n+2} = s\,a_{n+1} - p\,a_n$: jadi pemangkatannya menjadi
pengulangan. (iii) Pertanyaan 6 merupakan kasus $n = 2$ atas
\emph{teorema Cayley--Hamilton}, yang sah pada segala dimensi dan
dibuktikan pada jilid Tahun~2. (iv) Adapun matriks pendampingnya menutup
lingkarannya: karena setiap \omterm{pb:b1:matrices:1}{rekurensi linear} \emph{adalah} pangkat matriks, dengan
\omterm{def:b1:poly:def}{polinomial} $D$ yang sama muncul sebagai data trace-dan-determinan,
sehingga kalkulus sisanya memecahkan rekurensi dan menghitung pangkat
dalam satu gerakan.

\end{solution}
